\documentclass[12pt]{article}
\usepackage[utf8]{inputenc}
\usepackage[T1]{fontenc}
\usepackage{hyperref}
\title{Tontemboan Language Resources for K-12 Curriculum Development}
\author{Compiled by Claude Sonnet 5 (claude-sonnet-5) for Andr\'as Kornai}
\date{}
\begin{document}
\maketitle

\section{Introduction}
Tontemboan (ISO 639-3 tnt) is an Austronesian language of the North Minahasan branch, spoken in North Sulawesi, Indonesia. Omniglot gives about 150,000 speakers; the Wikipedia figure dates from 1990, so current numbers may be lower. The language is written in the Latin alphabet, and several spelling conventions appear to be in use. A 2024 study reports declining use among young people. Among surveyed high school students, only 3.4\% actively used Tontemboan with their family. The search found no school primers or textbooks. The searches turned up only a few academic and reference items.

\section{Primary Educational Resources}
No primers, textbooks or school curricula in Tontemboan turned up in the search results. This gap is itself a finding.

\section{Linguistic and Reference Materials}
\begin{itemize}
\item Wikipedia, ``Tontemboan language'': \url{https://en.wikipedia.org/wiki/Tontemboan_language}. The search summary says the classic reference works are the Dutch-language works of J.\ A.\ T.\ Schwarz from 1907--1908: a Tontemboansch-Nederlandsch dictionary, a grammar (\emph{Hoofdstukken uit de Spraakkunst van het Tontemboansch}) and three volumes of \emph{Tontemboansche Teksten}. I did not see digitized copies and give no URL for them.
\item ``Reconstruction of the Tontemboan-Indonesian Dictionary: Makela'i and Matana'i'' (2024, Indonesian-language article): \url{https://ejournal.unibabwi.ac.id/index.php/santhet/article/view/3287}. Download: \url{https://ejournal.unibabwi.ac.id/index.php/santhet/article/download/3287/2771}.
\item ``Prefixes In Tontemboan Language'' (morphology study): \url{https://journal.trunojoyo.ac.id/ijcst/article/view/8215}.
\item Related download from the same search: \url{https://ejournal.unibabwi.ac.id/index.php/santhet/article/download/4688/2781}. I did not confirm its contents from the snippet.
\item Omniglot overview with writing-system notes: \url{https://omniglot.com/writing/tontemboan.htm}.
\end{itemize}

\section{Religious and Community Resources}
Nothing specific to Tontemboan (scripture translations or community organizations) turned up in the search results.

\section{Digital and Audio Resources}
\begin{itemize}
\item Verbix page with sample verbs: \url{https://docs.verbix.com/Languages/Tontemboan}.
\end{itemize}
No audio or video collections turned up.

\section{Curriculum Development Context}
Speakers use Tontemboan mainly as a first language among adults, and children are shifting to Indonesian. Any curriculum would have to start from scratch: standardize the orthography, convert the Dutch-era dictionary and grammar and the 2024 dictionary into usable form, and collect new texts and audio. The most useful partners are probably the university researchers who authored the recent studies. I found no evidence of existing school provision.

\end{document}